\chapter{二阶及高阶微分方程}
前两章讨论的都是一阶微分方程. 本章讨论二阶及高阶微分方程，沿用第 1 章的记号，
用 $x=x(t)$ 表示未知函数，$t$ 表示自变量. 与一阶方程不同，高阶方程除少数可降阶
的类型外，一般不能通过初等积分求出通解. 因此本章先介绍几类可降阶的方程，再建立
线性方程的基本理论，最后讨论常系数线性方程的解法.

\section{可降阶的高阶方程}\label{sec:3.1}
本节介绍三种可降阶的方程类型及其解法. $n$ 阶微分方程的一般形式为
\begin{equation}\label{equ:3.1}
F(t,x,x',\dots,x^{(n)})=0.
\end{equation}
当 $n\ge2$ 时，称 (\ref{equ:3.1}) 为\textbf{高阶微分方程}. 与一阶方程的通解含有
一个任意常数类似，$n$ 阶方程 (\ref{equ:3.1}) 的通解一般含有 $n$ 个独立的任意常数.

\subsection{不显含未知函数 $x$ 的方程}
不显含未知函数 $x$，或更一般地，不显含未知函数及其直到 $k-1\ (k\ge1)$ 阶导数
的方程形如
\begin{equation}\label{equ:3.2}
F(t,x^{(k)},x^{(k+1)},\dots,x^{(n)})=0.
\end{equation}
令 $y=x^{(k)}$，则方程 (\ref{equ:3.2}) 化为关于 $y$ 的 $n-k$ 阶方程
\begin{equation}\label{equ:3.3}
F(t,y,y',\dots,y^{(n-k)})=0.
\end{equation}
若求得 (\ref{equ:3.3}) 的通解
\[
y=\varphi(t,c_1,c_2,\dots,c_{n-k}),
\]
即 $x^{(k)}=\varphi(t,c_1,c_2,\dots,c_{n-k})$，再对上式积分 $k$ 次，便得到方程
(\ref{equ:3.2}) 的通解，其中又增加了 $k$ 个任意常数.

\subsection{不显含自变量 $t$ 的方程}
不显含自变量 $t$ 的方程形如
\begin{equation}\label{equ:3.4}
F(x,x',\dots,x^{(n)})=0.
\end{equation}
此时取 $y=x'$ 作为新的未知函数，而把 $x$ 看作新的自变量. 由复合函数求导法则，
\[
\frac{dx}{dt}=y,\qquad
\frac{d^2x}{dt^2}=\frac{dy}{dt}=\frac{dy}{dx}\frac{dx}{dt}=y\frac{dy}{dx},
\]
\[
\frac{d^3x}{dt^3}=\frac{d}{dt}\left(y\frac{dy}{dx}\right)
=y\left(\frac{dy}{dx}\right)^2+y^2\frac{d^2y}{dx^2}.
\]
一般地，用数学归纳法可以证明，$x^{(k)}$ 可用
$y,\dfrac{dy}{dx},\dots,\dfrac{d^{k-1}y}{dx^{k-1}}\ (k\le n)$ 表示. 将这些表达式
代入 (\ref{equ:3.4})，得
\[
F\left(x,y,y\frac{dy}{dx},\,y\left(\frac{dy}{dx}\right)^2
+y^2\frac{d^2y}{dx^2},\dots\right)=0,
\]
即得到关于 $y$ 的 $n-1$ 阶方程
\begin{equation}\label{equ:3.5}
G\left(x,y,\frac{dy}{dx},\dots,\frac{d^{n-1}y}{dx^{n-1}}\right)=0.
\end{equation}
它比原方程 (\ref{equ:3.4}) 降低了一阶.

\subsection{全微分方程和积分因子}
若方程 (\ref{equ:3.1}) 的左端是某个 $n-1$ 阶微分表达式
$\Phi\left(t,x,\dfrac{dx}{dt},\dots,\dfrac{d^{n-1}x}{dt^{n-1}}\right)$ 对 $t$ 的全
导数，即
\begin{equation}\label{equ:3.6}
F\left(t,x,\frac{dx}{dt},\dots,\frac{d^n x}{dt^n}\right)
=\frac{d}{dt}\Phi\left(t,x,\frac{dx}{dt},\dots,\frac{d^{n-1}x}{dt^{n-1}}\right),
\end{equation}
则与一阶方程类似，称 (\ref{equ:3.1}) 为\textbf{全微分方程}. 此时显然有
\begin{equation}\label{equ:3.7}
\Phi\left(t,x,\frac{dx}{dt},\dots,\frac{d^{n-1}x}{dt^{n-1}}\right)=c_1.
\end{equation}
方程 (\ref{equ:3.7}) 是 $n-1$ 阶的. 若能求出 (\ref{equ:3.7}) 的通解
$x=\varphi(t,c_1,c_2,\dots,c_n)$，则它一定也是原方程 (\ref{equ:3.1}) 的通解.

有时 (\ref{equ:3.1}) 本身不是全微分方程，但乘以适当的因子
$\mu\left(t,x,\dfrac{dx}{dt},\dots,\dfrac{d^{n-1}x}{dt^{n-1}}\right)$ 后能成为全
微分方程，则称 $\mu$ 为方程 (\ref{equ:3.1}) 的\textbf{积分因子}.

\section{线性微分方程的基本理论}\label{sec:3.2}
线性微分方程是常微分方程中一类很重要的方程，它的理论发展得十分完善. 本节介绍
它的基本理论.

\subsection{线性微分方程的有关概念}
未知函数 $x$ 及其各阶导数 $\dfrac{dx}{dt},\dots,\dfrac{d^n x}{dt^n}$ 均为一次的
$n$ 阶微分方程称为 $n$ 阶\textbf{线性微分方程}，它的一般形式是
\begin{equation}\label{equ:3.8}
\frac{d^n x}{dt^n}+a_1(t)\frac{d^{n-1}x}{dt^{n-1}}+\cdots
+a_{n-1}(t)\frac{dx}{dt}+a_n(t)x=f(t),
\end{equation}
其中 $a_i(t)\ (i=1,2,\dots,n)$ 及 $f(t)$ 都是区间 $(a,b)$ 上的连续函数.

如果 $f(t)\equiv0$，则 (\ref{equ:3.8}) 变为
\begin{equation}\label{equ:3.9}
\frac{d^n x}{dt^n}+a_1(t)\frac{d^{n-1}x}{dt^{n-1}}+\cdots
+a_{n-1}(t)\frac{dx}{dt}+a_n(t)x=0,
\end{equation}
称 (\ref{equ:3.9}) 为 $n$ 阶\textbf{齐次线性方程}；相应地，称 (\ref{equ:3.8}) 为
$n$ 阶\textbf{非齐次线性方程}，并称 (\ref{equ:3.9}) 为与 (\ref{equ:3.8}) 对应的
齐次方程.

与一阶方程一样，首先要讨论解的存在唯一性.
\begin{thm}\label{thm:3.2.1}
若 $a_i(t)\ (i=1,2,\dots,n)$ 及 $f(t)$ 在区间 $(a,b)$ 上连续，则对任意
$t_0\in(a,b)$ 及任意 $x_0,x_0^{(1)},\dots,x_0^{(n-1)}$，方程 (\ref{equ:3.8})
存在唯一的解 $x=\varphi(t)$，满足初始条件
\[
\varphi(t_0)=x_0,\quad \varphi'(t_0)=x_0^{(1)},\quad\dots,\quad
\varphi^{(n-1)}(t_0)=x_0^{(n-1)}.
\]
\end{thm}
\begin{pf1}
将方程 (\ref{equ:3.8}) 化为一阶线性方程组，重复第 1 章定理 \ref{thm:1.1.1} 的
论证即可，此处从略.
\end{pf1}

如无特别声明，本节总假设 (\ref{equ:3.8}) 的系数 $a_i(t)\ (i=1,2,\dots,n)$ 及
$f(t)$ 在区间 $(a,b)$ 上连续.

为了叙述方便，引入记号
\begin{equation}\label{equ:3.10}
L[x]\equiv\frac{d^n x}{dt^n}+a_1(t)\frac{d^{n-1}x}{dt^{n-1}}+\cdots
+a_{n-1}(t)\frac{dx}{dt}+a_n(t)x,
\end{equation}
并称 $L$ 为\textbf{线性微分算子}. 当 $L$ 作用于函数 $x$ 时，就是指对 $x$ 施行
(\ref{equ:3.10}) 右端的运算. 例如，取 $x(t)=e^{\lambda t}$，则
\[
L[e^{\lambda t}]=\left[\lambda^n+a_1(t)\lambda^{n-1}
+a_2(t)\lambda^{n-2}+\cdots+a_{n-1}(t)\lambda+a_n(t)\right]e^{\lambda t}.
\]
于是 (\ref{equ:3.8}) 与 (\ref{equ:3.9}) 可分别简写为 $L[x]=f(t)$ 与 $L[x]=0$.

线性微分算子 $L$ 具有下面的性质.
\begin{prop}
\begin{description}
\item[(i)] $L[cx]=cL[x]$，其中 $c$ 为常数；
\item[(ii)] $L[x_1+x_2]=L[x_1]+L[x_2]$.
\end{description}
\end{prop}

\subsection{齐次线性方程解的性质和结构}
本节所讨论的齐次线性方程为 (\ref{equ:3.9})，总假设其系数在 $(a,b)$ 上连续.
\begin{thm}\label{thm:3.2.2}
若 $x_1(t),x_2(t),\dots,x_k(t)$ 是方程 (\ref{equ:3.9}) 的 $k$ 个解，则它们的
线性组合 $c_1x_1(t)+c_2x_2(t)+\cdots+c_kx_k(t)$ 也是 (\ref{equ:3.9}) 的解，其中
$c_1,c_2,\dots,c_k$ 为常数.
\end{thm}
\begin{pf1}
因为 $x_i(t)\ (i=1,2,\dots,k)$ 是 (\ref{equ:3.9}) 的解，故 $L[x_i(t)]=0$. 由
线性微分算子的性质可得
\[
L\left[c_1x_1(t)+c_2x_2(t)+\cdots+c_kx_k(t)\right]=0,
\]
即 $c_1x_1(t)+c_2x_2(t)+\cdots+c_kx_k(t)$ 是 (\ref{equ:3.9}) 的解.
\end{pf1}

若 $k=n$，即 $n$ 阶方程 (\ref{equ:3.9}) 有 $n$ 个解 $x_1(t),x_2(t),\dots,x_n(t)$，
则由定理 \ref{thm:3.2.2} 知，$\sum\limits_{i=1}^n c_ix_i(t)$ 也是 (\ref{equ:3.9})
的解，它含有 $n$ 个任意常数. 反过来，如果 (\ref{equ:3.9}) 的任一解 $\varphi(t)$
都可以表示为 $\varphi(t)=\sum\limits_{i=1}^n c_ix_i(t)$，则称
$x_1(t),x_2(t),\dots,x_n(t)$ 为 (\ref{equ:3.9}) 的\textbf{基本解组}. 自然要问：
$n$ 个解在什么条件下构成基本解组？或者说，
$\sum\limits_{i=1}^n c_ix_i(t)$ 在什么条件下是 (\ref{equ:3.9}) 的通解？为回答
这些问题，先介绍函数组线性相关、线性无关以及 Wronskian 行列式等概念.

\begin{deff}
设 $x_1(t),x_2(t),\dots,x_k(t)$ 是定义在区间 $(a,b)$ 上的 $k$ 个函数. 如果存在
不全为零的常数 $c_1,c_2,\dots,c_k$，使得
\[
c_1x_1(t)+c_2x_2(t)+\cdots+c_kx_k(t)=0
\]
在 $(a,b)$ 上恒成立，则称这些函数在 $(a,b)$ 上\textbf{线性相关}；否则，称它们
在 $(a,b)$ 上\textbf{线性无关}.
\end{deff}

\begin{rmk}
函数 $1,t,t^2,\dots,t^n$ 在任何区间上都是线性无关的. 事实上，恒等式
\[
c_0+c_1t+c_2t^2+\cdots+c_nt^n=0
\]
只有当 $c_i=0\ (i=0,1,\dots,n)$ 时才成立. 若某个 $c_i\neq0$，则左端是一个不高于
$n$ 次的多项式，它最多有 $n$ 个不同的根，不可能在整个区间上恒为零.
\end{rmk}

\begin{rmk}
在函数组 $x_1(t),x_2(t),\dots,x_k(t)$ 中，若有一个函数在 $(a,b)$ 上恒等于零，
则该函数组在 $(a,b)$ 上线性相关.
\end{rmk}

\begin{rmk}
考虑由两个函数构成的函数组 $x_1(t),x_2(t)$. 如果 $\dfrac{x_1(t)}{x_2(t)}$
（或 $\dfrac{x_2(t)}{x_1(t)}$）在 $(a,b)$ 上有定义，则它们线性无关的充要条件是
该比值在 $(a,b)$ 上不恒为常数. 例如，$\sin t$ 与 $\cos t$ 在任何区间上都线性
无关，而 $\cos^2t$ 与 $\sin^2t-1$ 在任何区间上都线性相关. 此外，线性相关与
线性无关还依赖于所取的区间：函数 $x_1(t)=|t|$ 与 $x_2(t)=t$ 在
$(-\infty,+\infty)$ 上线性无关，但在 $(-\infty,0)$ 与 $(0,+\infty)$ 上都线性相关，
因为
\[
\frac{x_1(t)}{x_2(t)}=
\begin{cases}
-1, & t<0,\\
1, & t>0
\end{cases}
\]
在 $(-\infty,+\infty)$ 上不是常数，而在 $(-\infty,0)$ 与 $(0,+\infty)$ 上都是常数.
\end{rmk}

下面建立线性相关与线性无关的判别法则. 设 $x_1(t),x_2(t),\dots,x_k(t)$ 是定义
在 $(a,b)$ 上的 $k$ 个 $k-1$ 次可微函数，由它们构成的行列式
\[
W[x_1(t),x_2(t),\dots,x_k(t)]\equiv
\begin{vmatrix}
x_1(t) & x_2(t) & \dots & x_k(t)\\
x_1'(t) & x_2'(t) & \dots & x_k'(t)\\
\vdots & \vdots & & \vdots\\
x_1^{(k-1)}(t) & x_2^{(k-1)}(t) & \dots & x_k^{(k-1)}(t)
\end{vmatrix},
\]
称为这些函数的 \textbf{Wronskian 行列式}，也记作 $W(t)$.

\begin{thm}\label{thm:3.2.3}
如果函数组 $x_1(t),x_2(t),\dots,x_n(t)$ 在 $(a,b)$ 上线性相关，则它们的
Wronskian 行列式在 $(a,b)$ 上恒等于零.
\end{thm}
\begin{pf1}
由假设，存在不全为零的常数 $c_1,c_2,\dots,c_n$，使得
\[
c_1x_1(t)+c_2x_2(t)+\cdots+c_nx_n(t)\equiv0,\quad t\in(a,b).
\]
依次对 $t$ 求导 $n-1$ 次，得到 $n$ 个恒等式
\[
\begin{cases}
c_1x_1(t)+c_2x_2(t)+\cdots+c_nx_n(t)=0,\\
c_1x_1'(t)+c_2x_2'(t)+\cdots+c_nx_n'(t)=0,\\
\quad\vdots\\
c_1x_1^{(n-1)}(t)+c_2x_2^{(n-1)}(t)+\cdots+c_nx_n^{(n-1)}(t)=0.
\end{cases}
\]
这是关于 $c_1,c_2,\dots,c_n$ 的齐次线性方程组，其系数行列式恰为 $W(t)$. 由于它
有非零解，系数行列式必为零，即 $W(t)\equiv0\ (t\in(a,b))$.
\end{pf1}

\begin{cor}\label{cor:3.2.1}
如果函数组 $x_1(t),x_2(t),\dots,x_n(t)$ 的 Wronskian 行列式在 $(a,b)$ 上某点
$t_0$ 处不等于零，即 $W(t_0)\neq0$，则该函数组在 $(a,b)$ 上线性无关.
\end{cor}

\begin{rmk}
定理 \ref{thm:3.2.3} 的逆命题不一定成立，即 $W(t)\equiv0$ 只是函数组线性相关的
必要条件，并非充分条件. 例如，函数
\[
x_1(t)=
\begin{cases}
t^2, & t\ge0,\\
0, & t<0
\end{cases}
\quad\text{与}\quad
x_2(t)=
\begin{cases}
0, & t\ge0,\\
t^2, & t<0
\end{cases}
\]
满足 $W[x_1(t),x_2(t)]\equiv0$，但它们在 $(-\infty,+\infty)$ 上线性无关. 事实上，
若 $c_1x_1(t)+c_2x_2(t)\equiv0$，则当 $t<0$ 时推得 $c_2=0$，当 $t\ge0$ 时推得
$c_1=0$，故 $x_1(t),x_2(t)$ 线性无关.

应当指出，如果 $x_1(t),x_2(t),\dots,x_n(t)$ 是齐次方程 (\ref{equ:3.9}) 的
$n$ 个解，那么 $W(t)\equiv0$ 就成为该解组线性相关的充要条件，这可由下面的定理
推出.
\end{rmk}

\begin{thm}\label{thm:3.2.4}
如果 $x_1(t),x_2(t),\dots,x_n(t)$ 是方程 (\ref{equ:3.9}) 在 $(a,b)$ 上的
$n$ 个线性无关的解，则它们的 Wronskian 行列式在 $(a,b)$ 上任何点都不等于零，
即 $W(t)\neq0\ (t\in(a,b))$.
\end{thm}
\begin{pf1}
采用反证法. 设有 $t_0\in(a,b)$ 使 $W(t_0)=0$. 考虑关于 $c_1,c_2,\dots,c_n$ 的
齐次线性方程组
\[
\begin{cases}
c_1x_1(t_0)+c_2x_2(t_0)+\cdots+c_nx_n(t_0)=0,\\
c_1x_1'(t_0)+c_2x_2'(t_0)+\cdots+c_nx_n'(t_0)=0,\\
\quad\vdots\\
c_1x_1^{(n-1)}(t_0)+c_2x_2^{(n-1)}(t_0)+\cdots+c_nx_n^{(n-1)}(t_0)=0.
\end{cases}
\]
其系数行列式为 $W(t_0)=0$，故它有非零解 $c_1,c_2,\dots,c_n$. 以这组解构造函数
\[
x(t)=c_1x_1(t)+c_2x_2(t)+\cdots+c_nx_n(t),\quad a<t<b,
\]
由定理 \ref{thm:3.2.2} 知 $x(t)$ 是 (\ref{equ:3.9}) 的解. 又由上述方程组知
$x(t)$ 满足零初始条件
\[
x(t_0)=x'(t_0)=\cdots=x^{(n-1)}(t_0)=0.
\]
但 $x(t)\equiv0$ 也是 (\ref{equ:3.9}) 满足该初始条件的解，由解的唯一性得
$x(t)\equiv0\ (a<t<b)$，即
\[
c_1x_1(t)+c_2x_2(t)+\cdots+c_nx_n(t)\equiv0.
\]
这与 $x_1(t),x_2(t),\dots,x_n(t)$ 线性无关矛盾.
\end{pf1}

\begin{cor}\label{cor:3.2.2}
设 $x_1(t),x_2(t),\dots,x_n(t)$ 是方程 (\ref{equ:3.9}) 在 $(a,b)$ 上的 $n$ 个解.
若存在 $t_0\in(a,b)$ 使 $W(t_0)=0$，则该解组在 $(a,b)$ 上线性相关.
\end{cor}

\begin{cor}\label{cor:3.2.3}
方程 (\ref{equ:3.9}) 的 $n$ 个解 $x_1(t),x_2(t),\dots,x_n(t)$ 在 $(a,b)$ 上线性
无关的充要条件是：存在 $t_0\in(a,b)$，使得 $W(t_0)\neq0$.
\end{cor}

由上述讨论可知，由 $n$ 阶齐次线性方程 (\ref{equ:3.9}) 的 $n$ 个解构成的
Wronskian 行列式，或者在 $(a,b)$ 上恒等于零，或者在 $(a,b)$ 上处处不等于零.

由定理 \ref{thm:3.2.1} 知，方程 (\ref{equ:3.9}) 满足初始条件
\[
\begin{cases}
x_1(t_0)=1,\ x_1'(t_0)=0,\ \dots,\ x_1^{(n-1)}(t_0)=0,\\
x_2(t_0)=0,\ x_2'(t_0)=1,\ \dots,\ x_2^{(n-1)}(t_0)=0,\\
\quad\vdots\\
x_n(t_0)=0,\ x_n'(t_0)=0,\ \dots,\ x_n^{(n-1)}(t_0)=1
\end{cases}
\]
的解 $x_1(t),x_2(t),\dots,x_n(t)$ 一定存在. 由于
$W(t_0)=1\neq0$，由推论 \ref{cor:3.2.3} 知这 $n$ 个解线性无关，于是有

\begin{thm}\label{thm:3.2.5}
$n$ 阶齐次线性方程 (\ref{equ:3.9}) 一定存在 $n$ 个线性无关的解.
\end{thm}

下面的定理给出了线性无关解组、基本解组与通解之间的关系.
\begin{thm}\label{thm:3.2.6}
如果 $x_1(t),x_2(t),\dots,x_n(t)$ 是 $n$ 阶齐次方程 (\ref{equ:3.9}) 的
$n$ 个线性无关的解，则它一定是该方程的基本解组，即 (\ref{equ:3.9}) 的任一解
$x(t)$ 都可表示为
\[
x(t)=c_1x_1(t)+c_2x_2(t)+\cdots+c_nx_n(t).
\]
\end{thm}
\begin{pf1}
设 $x(t)$ 是 (\ref{equ:3.9}) 的任一解，且满足
\[
x(t_0)=x_0,\quad x'(t_0)=x_0',\quad\dots,\quad
x^{(n-1)}(t_0)=x_0^{(n-1)}.
\]
考虑关于 $c_1,c_2,\dots,c_n$ 的方程组
\[
\begin{cases}
x_0=c_1x_1(t_0)+c_2x_2(t_0)+\cdots+c_nx_n(t_0),\\
x_0'=c_1x_1'(t_0)+c_2x_2'(t_0)+\cdots+c_nx_n'(t_0),\\
\quad\vdots\\
x_0^{(n-1)}=c_1x_1^{(n-1)}(t_0)+c_2x_2^{(n-1)}(t_0)
+\cdots+c_nx_n^{(n-1)}(t_0).
\end{cases}
\]
其系数行列式是 $W(t_0)\neq0$，故方程组有唯一解 $c_1,c_2,\dots,c_n$. 令
$\varphi(t)=\sum\limits_{i=1}^n c_ix_i(t)$，由定理 \ref{thm:3.2.2} 知
$\varphi(t)$ 是 (\ref{equ:3.9}) 的解，且与 $x(t)$ 满足相同的初始条件. 由解的
唯一性得 $\varphi(t)=x(t)$，即 $x(t)=\sum\limits_{i=1}^n c_ix_i(t)$.
\end{pf1}

由定理 \ref{thm:3.2.6} 立即得到齐次线性方程的通解结构.
\begin{thm}\label{thm:3.2.7}
如果 $x_1(t),x_2(t),\dots,x_n(t)$ 是方程 (\ref{equ:3.9}) 的 $n$ 个线性无关解，
则方程 (\ref{equ:3.9}) 的通解为
\begin{equation}\label{equ:3.11}
x(t)=c_1x_1(t)+c_2x_2(t)+\cdots+c_nx_n(t),
\end{equation}
其中 $c_1,c_2,\dots,c_n$ 为任意常数.
\end{thm}

综合上面的讨论，得到下列等价结论.
\begin{thm}\label{thm:3.2.8}
设 $x_1(t),x_2(t),\dots,x_n(t)$ 是方程 (\ref{equ:3.9}) 的 $n$ 个解，则下列命题
等价：
\begin{enumerate}
\item 方程 (\ref{equ:3.9}) 的通解为
$\displaystyle x(t)=\sum_{i=1}^n c_ix_i(t)$；
\item $x_1(t),x_2(t),\dots,x_n(t)$ 是 (\ref{equ:3.9}) 的基本解组；
\item $x_1(t),x_2(t),\dots,x_n(t)$ 在 $(a,b)$ 上线性无关；
\item 在 $(a,b)$ 上存在一点 $t_0$，使得 $W(t_0)\neq0$；
\item 在 $(a,b)$ 上任一点 $t$，都有 $W(t)\neq0$.
\end{enumerate}
\end{thm}

Wronskian 行列式还满足下面的 Liouville 公式，它在具体计算中十分有用.
\begin{thm}\label{thm:3.2.9}
设 $a_1(t),a_2(t),\dots,a_n(t)$ 在区间 $I$ 上连续，
$x_1(t),x_2(t),\dots,x_n(t)$ 是 $n$ 阶齐次线性方程
\[
x^{(n)}+a_1(t)x^{(n-1)}+a_2(t)x^{(n-2)}+\cdots+a_n(t)x=0
\]
在 $I$ 上的 $n$ 个解，$W(t)=W[x_1,x_2,\dots,x_n](t)$ 为这组解的 Wronskian
行列式，则对任意 $t_0\in I$，有
\[
W(t)=W(t_0)\exp\left(-\int_{t_0}^{t}a_1(s)\,ds\right),\quad t\in I.
\]
\end{thm}
\begin{pf1}
行列式对自变量求导遵循逐行求导法则：行列式的导数等于分别对每一行求导、其余行
保持不变后所得各行列式之和，即
\[
W'(t)=\sum_{k=1}^{n}
\begin{vmatrix}
x_1 & x_2 & \cdots & x_n\\
x_1' & x_2' & \cdots & x_n'\\
\vdots & \vdots & & \vdots\\
\left(x_1^{(k-1)}\right)' & \left(x_2^{(k-1)}\right)' & \cdots
& \left(x_n^{(k-1)}\right)'\\
\vdots & \vdots & & \vdots\\
x_1^{(n-1)} & x_2^{(n-1)} & \cdots & x_n^{(n-1)}
\end{vmatrix}.
\]
当 $k=1,2,\dots,n-1$ 时，被求导的第 $k$ 行与下一行完全相同，相应的行列式为零，
故求和式中只保留 $k=n$ 这一项：
\[
W'(t)=
\begin{vmatrix}
x_1 & x_2 & \cdots & x_n\\
x_1' & x_2' & \cdots & x_n'\\
\vdots & \vdots & & \vdots\\
x_1^{(n-2)} & x_2^{(n-2)} & \cdots & x_n^{(n-2)}\\
x_1^{(n)} & x_2^{(n)} & \cdots & x_n^{(n)}
\end{vmatrix}.
\]
由于 $x_i(t)$ 是齐次方程的解，故
\[
x_i^{(n)}=-a_1(t)x_i^{(n-1)}-a_2(t)x_i^{(n-2)}-\cdots-a_n(t)x_i,
\quad i=1,2,\dots,n.
\]
将上式代入最后一行，利用行列式对行的线性性质拆分，则 $a_j\ (j\ge2)$ 对应的
各项均出现两行相同的情形，故全为零，只剩 $a_1$ 对应的项：
\[
W'(t)=-a_1(t)
\begin{vmatrix}
x_1 & x_2 & \cdots & x_n\\
x_1' & x_2' & \cdots & x_n'\\
\vdots & \vdots & & \vdots\\
x_1^{(n-2)} & x_2^{(n-2)} & \cdots & x_n^{(n-2)}\\
x_1^{(n-1)} & x_2^{(n-1)} & \cdots & x_n^{(n-1)}
\end{vmatrix}
=-a_1(t)W(t).
\]
于是 $W(t)$ 满足一阶线性齐次方程
\[
W'(t)+a_1(t)W(t)=0.
\]
分离变量得
\[
\frac{dW}{W}=-a_1(t)\,dt,
\]
在 $I$ 上从 $t_0$ 积分到 $t$，得
\[
\int_{W(t_0)}^{W(t)}\frac{du}{u}
=-\int_{t_0}^{t}a_1(s)\,ds.
\]
两端取指数，即得
\[
W(t)=W(t_0)\exp\left(-\int_{t_0}^{t}a_1(s)\,ds\right).
\]
\end{pf1}

\begin{rmk}
由 Liouville 公式可知，方程 (\ref{equ:3.9}) 的解组的 Wronskian 行列式 $W(t)$
若在 $(a,b)$ 内某点为零，则在整个 $(a,b)$ 上恒为零；若在某点不为零，则在整个
$(a,b)$ 上恒不为零.
\end{rmk}

\begin{rmk}
对二阶齐次线性方程 $x''+p(t)x'+q(t)x=0$，若已知它的一个非零解 $x_1(t)$，则可
由 Liouville 公式求出与 $x_1(t)$ 线性无关的另一个解，从而写出通解.

事实上，设 $x(t)$ 是另一个解，则 $W[x_1,x]=x_1x'-x_1'x$ 满足
\[
x_1x'-x_1'x=c\exp\left(-\int p(t)\,dt\right).
\]
上式两端乘以 $\dfrac{1}{x_1^2}$，得
\[
\frac{d}{dt}\left(\frac{x}{x_1}\right)
=\frac{c}{x_1^2}\exp\left(-\int p(t)\,dt\right),
\]
积分得
\[
\frac{x}{x_1}=c\int\frac{1}{x_1^2}\exp\left(-\int p(t)\,dt\right)dt+c_1.
\]
取 $c=1,c_1=0$，则
\[
x_2(t)=x_1(t)\int\frac{1}{x_1^2(t)}
\exp\left(-\int p(t)\,dt\right)dt
\]
就是与 $x_1(t)$ 线性无关的另一个解，原方程的通解为 $x=c_1x_1+c_2x_2$.
\end{rmk}

\subsection{非齐次线性方程解的结构}
本节讨论非齐次线性方程 (\ref{equ:3.8}). 与一阶线性方程一样，它的通解由对应的
齐次方程的通解与非齐次方程本身的一个特解构成.

\begin{thm}\label{thm:3.2.10}
$n$ 阶非齐次线性方程 (\ref{equ:3.8}) 的通解等于它所对应的齐次方程
(\ref{equ:3.9}) 的通解与 (\ref{equ:3.8}) 本身的一个特解之和.
\end{thm}
\begin{pf1}
设 $x^*$ 是 (\ref{equ:3.8}) 的一个特解，$\tilde{x}$ 是 (\ref{equ:3.9}) 的通解.
首先，由于
\[
L[x^*+\tilde{x}]=L[x^*]+L[\tilde{x}]=f(t)+0=f(t),
\]
故 $x=x^*+\tilde{x}$ 是 (\ref{equ:3.8}) 的解.

其次，设 $\bar{x}$ 是 (\ref{equ:3.8}) 的任一解，则
\[
L[\bar{x}-x^*]=L[\bar{x}]-L[x^*]=f(t)-f(t)=0,
\]
故 $\bar{x}-x^*$ 是 (\ref{equ:3.9}) 的解，记为 $\tilde{x}_0$，其中 $\tilde{x}_0$ 是
由 $\tilde{x}$ 中的任意常数取某组确定值得到的. 于是 $\bar{x}=x^*+\tilde{x}_0$，
这说明 $x^*+\tilde{x}$ 包含了 (\ref{equ:3.8}) 的所有解.
\end{pf1}

\begin{thm}\label{thm:3.2.11}
设 $x_1(t),x_2(t)$ 分别是非齐次线性方程 $L[x]=f_1(t)$ 与 $L[x]=f_2(t)$ 的解，
则 $x_1(t)+x_2(t)$ 是方程 $L[x]=f_1(t)+f_2(t)$ 的解.
\end{thm}
\begin{pf1}
由假设 $L[x_1]=f_1,\ L[x_2]=f_2$，故
\[
L[x_1+x_2]=L[x_1]+L[x_2]=f_1(t)+f_2(t),
\]
即 $x_1(t)+x_2(t)$ 是 $L[x]=f_1(t)+f_2(t)$ 的解.
\end{pf1}

\subsection{常数变易法}
与一阶线性方程类似，$n$ 阶非齐次线性方程也可以利用对应齐次方程的 $n$ 个线性
无关解求其特解，这种方法称为\textbf{常数变易法}.

设 (\ref{equ:3.11}) 是齐次方程 (\ref{equ:3.9}) 的通解. 为求 (\ref{equ:3.8})
的特解，把其中的常数 $c_i$ 换为待定函数 $c_i(t)$，即令
\begin{equation}\label{equ:3.12}
x=c_1(t)x_1(t)+c_2(t)x_2(t)+\cdots+c_n(t)x_n(t).
\end{equation}
将 (\ref{equ:3.12}) 代入 (\ref{equ:3.8})，只能得到 $c_i(t)$ 满足的一个条件. 由于
待定函数有 $n$ 个，还需再给出 $n-1$ 个条件. 为运算方便，按下面的方式给出这些
条件.

对 (\ref{equ:3.12}) 求导，得
\[
\begin{aligned}
x'={}&c_1(t)x_1'(t)+c_2(t)x_2'(t)+\cdots+c_n(t)x_n'(t)\\
&+c_1'(t)x_1(t)+c_2'(t)x_2(t)+\cdots+c_n'(t)x_n(t).
\end{aligned}
\]
令
\[
c_1'(t)x_1(t)+c_2'(t)x_2(t)+\cdots+c_n'(t)x_n(t)=0,
\]
则有 $x'=c_1x_1'+\cdots+c_nx_n'$. 对上式继续求导，并同样令含有 $c_i'(t)$ 的部分
为零，又得到一个条件
\[
c_1'(t)x_1'(t)+c_2'(t)x_2'(t)+\cdots+c_n'(t)x_n'(t)=0
\]
和表达式 $x''=c_1x_1''+\cdots+c_nx_n''$. 如此继续，共得到 $n-1$ 个条件，最后
一个为
\[
c_1'(t)x_1^{(n-2)}(t)+c_2'(t)x_2^{(n-2)}(t)+\cdots
+c_n'(t)x_n^{(n-2)}(t)=0,
\]
此时 $x^{(n-1)}=c_1x_1^{(n-1)}+\cdots+c_nx_n^{(n-1)}$. 再对上式求导，得
\[
\begin{aligned}
x^{(n)}={}&c_1(t)x_1^{(n)}(t)+c_2(t)x_2^{(n)}(t)+\cdots+c_n(t)x_n^{(n)}(t)\\
&+c_1'(t)x_1^{(n-1)}(t)+c_2'(t)x_2^{(n-1)}(t)
+\cdots+c_n'(t)x_n^{(n-1)}(t).
\end{aligned}
\]
将上述 $x,x',\dots,x^{(n)}$ 代入 (\ref{equ:3.8})，并注意 $x_i$ 是 (\ref{equ:3.9})
的解，便得
\[
c_1'(t)x_1^{(n-1)}(t)+c_2'(t)x_2^{(n-1)}(t)+\cdots
+c_n'(t)x_n^{(n-1)}(t)=f(t).
\]
于是 $c_1'(t),\dots,c_n'(t)$ 满足线性方程组
\[
\begin{cases}
c_1'x_1+c_2'x_2+\cdots+c_n'x_n=0,\\
c_1'x_1'+c_2'x_2'+\cdots+c_n'x_n'=0,\\
\quad\vdots\\
c_1'x_1^{(n-2)}+c_2'x_2^{(n-2)}+\cdots+c_n'x_n^{(n-2)}=0,\\
c_1'x_1^{(n-1)}+c_2'x_2^{(n-1)}+\cdots+c_n'x_n^{(n-1)}=f(t).
\end{cases}
\]
其系数行列式恰为 $W[x_1,x_2,\dots,x_n]\neq0$，故解 $c_i'(t)=\varphi_i(t)\ (i=1,2,
\dots,n)$ 唯一确定. 积分得
\[
c_i(t)=\int\varphi_i(t)\,dt+\gamma_i,\quad i=1,2,\dots,n,
\]
其中 $\gamma_i$ 为任意常数. 代回 (\ref{equ:3.12})，即得 (\ref{equ:3.8}) 的通解
\[
x(t)=\sum_{i=1}^n\gamma_ix_i(t)
+\sum_{i=1}^n x_i(t)\int\varphi_i(t)\,dt.
\]
取 $\gamma_i=0\ (i=1,2,\dots,n)$，便得到 (\ref{equ:3.8}) 的一个特解
$x=\sum\limits_{i=1}^n x_i(t)\int\varphi_i(t)\,dt$.

\section{线性齐次常系数方程}\label{sec:3.3}
本节讨论系数为常数的齐次线性方程. 为此，先简要介绍复值函数与 Euler 公式.

\subsection{复值函数}
若 $\varphi(t)$ 与 $\psi(t)$ 是区间 $(a,b)$ 上的实函数，$\mathrm{i}=\sqrt{-1}$ 为
虚数单位，则称
\[
z(t)=\varphi(t)+\mathrm{i}\psi(t)
\]
为 $(a,b)$ 上的\textbf{复值函数}. 若 $\varphi(t),\psi(t)$ 在 $(a,b)$ 上连续（可微），
则称 $z(t)$ 在 $(a,b)$ 上连续（可微），并定义
\[
\frac{dz}{dt}=\frac{d\varphi}{dt}+\mathrm{i}\frac{d\psi}{dt}.
\]
由定义容易验证，若 $z_1(t),z_2(t)$ 可微，$c$ 为复常数，则
\[
\begin{aligned}
\frac{d}{dt}\left[z_1(t)+z_2(t)\right]&=\frac{dz_1}{dt}+\frac{dz_2}{dt},\\
\frac{d}{dt}\left[cz_1(t)\right]&=c\frac{dz_1}{dt},\\
\frac{d}{dt}\left[z_1(t)z_2(t)\right]&=\frac{dz_1}{dt}z_2(t)
+z_1(t)\frac{dz_2}{dt}.
\end{aligned}
\]
在讨论常系数线性方程时，复指数函数 $e^{\lambda t}\ (\lambda=\alpha+\mathrm{i}
\beta)$ 起着重要作用. 由指数函数的级数展开，得
\[
\begin{aligned}
e^{\mathrm{i}\beta t}
&=1+\mathrm{i}\beta t+\frac{(\mathrm{i}\beta t)^2}{2!}
+\frac{(\mathrm{i}\beta t)^3}{3!}+\cdots\\
&=\left[1-\frac{(\beta t)^2}{2!}+\frac{(\beta t)^4}{4!}-\cdots\right]
+\mathrm{i}\left[\beta t-\frac{(\beta t)^3}{3!}
+\frac{(\beta t)^5}{5!}-\cdots\right]\\
&=\cos\beta t+\mathrm{i}\sin\beta t.
\end{aligned}
\]
因此
\begin{equation}\label{equ:3.13}
e^{\mathrm{i}\beta t}=\cos\beta t+\mathrm{i}\sin\beta t,\qquad
e^{-\mathrm{i}\beta t}=\cos\beta t-\mathrm{i}\sin\beta t.
\end{equation}
由此还可解得
\begin{equation}\label{equ:3.14}
\cos\beta t=\frac12\left(e^{\mathrm{i}\beta t}+e^{-\mathrm{i}\beta t}\right),
\qquad
\sin\beta t=\frac{1}{2\mathrm{i}}\left(e^{\mathrm{i}\beta t}
-e^{-\mathrm{i}\beta t}\right).
\end{equation}
公式 (\ref{equ:3.13}) 与 (\ref{equ:3.14}) 统称为 \textbf{Euler 公式}.

记 $\bar\lambda=\alpha-\mathrm{i}\beta$ 为 $\lambda=\alpha+\mathrm{i}\beta$ 的
共轭，由 Euler 公式可得 $\overline{e^{\lambda t}}=e^{\bar\lambda t}$. 复指数
函数与实指数函数有类似的性质. 对复常数 $\lambda_1,\lambda_2,\lambda$，有
\[
e^{(\lambda_1+\lambda_2)t}=e^{\lambda_1t}e^{\lambda_2t},\qquad
\frac{d}{dt}e^{\lambda t}=\lambda e^{\lambda t}.
\]
后一式的验证如下：
\[
\begin{aligned}
\frac{d}{dt}e^{\lambda t}
&=\frac{d}{dt}\left[e^{\alpha t}(\cos\beta t
+\mathrm{i}\sin\beta t)\right]\\
&=\alpha e^{\alpha t}(\cos\beta t+\mathrm{i}\sin\beta t)
+e^{\alpha t}(-\beta\sin\beta t+\mathrm{i}\beta\cos\beta t)\\
&=e^{\alpha t}\left[\alpha(\cos\beta t+\mathrm{i}\sin\beta t)
+\mathrm{i}\beta(\mathrm{i}\sin\beta t+\cos\beta t)\right]\\
&=(\alpha+\mathrm{i}\beta)e^{\alpha t}(\cos\beta t
+\mathrm{i}\sin\beta t)=\lambda e^{\lambda t}.
\end{aligned}
\]

下面给出线性方程复值解的概念. 设方程 (\ref{equ:3.9}) 中的 $a_i(t)$ 为实函数，
若实变量复值函数 $x=z(t)=\varphi(t)+\mathrm{i}\psi(t)$ 满足该方程，则称 $z(t)$
为方程 (\ref{equ:3.9}) 的\textbf{复值解}.

\begin{thm}\label{thm:3.3.1}
若方程 (\ref{equ:3.9}) 的系数 $a_i(t)$ 均为实值函数，
$z(t)=\varphi(t)+\mathrm{i}\psi(t)$ 是 (\ref{equ:3.9}) 的复值解，则 $z(t)$ 的
实部 $\varphi(t)$、虚部 $\psi(t)$ 以及共轭 $\overline{z(t)}$ 都是 (\ref{equ:3.9})
的解.
\end{thm}
\begin{pf1}
由 $L$ 的性质，有
\[
L[z]=L\left[\varphi(t)+\mathrm{i}\psi(t)\right]
=L[\varphi(t)]+\mathrm{i}L[\psi(t)]=0.
\]
由于 $L[\varphi(t)]$ 与 $L[\psi(t)]$ 均为实值函数，故
\[
L[\varphi(t)]=L[\psi(t)]=0,
\]
即 $\varphi(t)$ 与 $\psi(t)$ 都是 (\ref{equ:3.9}) 的解. 又
\[
L\left[\overline{z(t)}\right]=\overline{L[z(t)]}
=L[\varphi(t)]-\mathrm{i}L[\psi(t)]=0,
\]
故 $\overline{z(t)}$ 也是 (\ref{equ:3.9}) 的解.
\end{pf1}

\subsection{常系数齐次线性方程}
考虑 $n$ 阶常系数齐次线性方程
\begin{equation}\label{equ:3.15}
L[x]=\frac{d^n x}{dt^n}+a_1\frac{d^{n-1}x}{dt^{n-1}}+\cdots
+a_{n-1}\frac{dx}{dt}+a_nx=0,
\end{equation}
其中 $a_1,a_2,\dots,a_n$ 为常数. 由 3.2 节的理论，求 (\ref{equ:3.15}) 的通解只
需求出它的基本解组. 下面介绍待定指数函数法.

尝试指数函数形式的解
\begin{equation}\label{equ:3.16}
x=e^{\lambda t}.
\end{equation}
将 (\ref{equ:3.16}) 代入 (\ref{equ:3.15})，得
\[
L[e^{\lambda t}]=\left(\lambda^n+a_1\lambda^{n-1}
+a_2\lambda^{n-2}+\cdots+a_{n-1}\lambda+a_n\right)e^{\lambda t}=0.
\]
因此，$e^{\lambda t}$ 是 (\ref{equ:3.15}) 的解的充要条件是 $\lambda$ 为代数方程
\begin{equation}\label{equ:3.17}
F(\lambda)\equiv\lambda^n+a_1\lambda^{n-1}
+a_2\lambda^{n-2}+\cdots+a_{n-1}\lambda+a_n=0
\end{equation}
的根. 方程 (\ref{equ:3.17}) 称为 (\ref{equ:3.15}) 的\textbf{特征方程}，它的根
称为 (\ref{equ:3.15}) 的\textbf{特征根}.

于是问题归结为求特征根. 下面按特征根的不同情形分别讨论.

\textbf{1. 特征根是单根的情形}

设 $\lambda_1,\lambda_2,\dots,\lambda_n$ 是 (\ref{equ:3.17}) 的 $n$ 个互不相同的
根，则相应得到 (\ref{equ:3.15}) 的 $n$ 个解
\begin{equation}\label{equ:3.18}
e^{\lambda_1t},\ e^{\lambda_2t},\ \dots,\ e^{\lambda_nt}.
\end{equation}

\begin{pf1}
考察 (\ref{equ:3.18}) 的 Wronskian 行列式：
\[
\begin{aligned}
W\left[e^{\lambda_1t},e^{\lambda_2t},\dots,e^{\lambda_nt}\right]
&=
\begin{vmatrix}
e^{\lambda_1t} & e^{\lambda_2t} & \cdots & e^{\lambda_nt}\\
\lambda_1e^{\lambda_1t} & \lambda_2e^{\lambda_2t} & \cdots
& \lambda_ne^{\lambda_nt}\\
\vdots & \vdots & & \vdots\\
\lambda_1^{n-1}e^{\lambda_1t} & \lambda_2^{n-1}e^{\lambda_2t} & \cdots
& \lambda_n^{n-1}e^{\lambda_nt}
\end{vmatrix}\\
&=e^{(\lambda_1+\lambda_2+\cdots+\lambda_n)t}
\begin{vmatrix}
1 & 1 & \cdots & 1\\
\lambda_1 & \lambda_2 & \cdots & \lambda_n\\
\vdots & \vdots & & \vdots\\
\lambda_1^{n-1} & \lambda_2^{n-1} & \cdots & \lambda_n^{n-1}
\end{vmatrix}.
\end{aligned}
\]
最后一个行列式是 Vandermonde 行列式，它等于
$\prod\limits_{1\le j<k\le n}(\lambda_k-\lambda_j)$. 由于 $\lambda_i$ 互不相同，
该行列式不为零，从而 $W\neq0$，故解组 (\ref{equ:3.18}) 线性无关.
\end{pf1}

若 $\lambda_1,\dots,\lambda_n$ 均为实数，则 (\ref{equ:3.18}) 是 (\ref{equ:3.15}) 的
$n$ 个线性无关实值解，通解为
\[
x(t)=c_1e^{\lambda_1t}+c_2e^{\lambda_2t}+\cdots+c_ne^{\lambda_nt},
\]
其中 $c_i$ 为任意常数.

若特征根中有复数，则由于系数为实常数，复根必成对共轭出现. 设
$\lambda_1=\alpha+\mathrm{i}\beta$ 为特征根，则
$\lambda_2=\alpha-\mathrm{i}\beta$ 也是特征根，对应的两个复值解为
\[
e^{(\alpha+\mathrm{i}\beta)t}=e^{\alpha t}(\cos\beta t+\mathrm{i}\sin\beta t),
\quad
e^{(\alpha-\mathrm{i}\beta)t}=e^{\alpha t}(\cos\beta t-\mathrm{i}\sin\beta t).
\]
由定理 \ref{thm:3.3.1}，它们的实部与虚部也是 (\ref{equ:3.15}) 的解，于是得到
两个线性无关的实值解
\[
e^{\alpha t}\cos\beta t,\qquad e^{\alpha t}\sin\beta t.
\]

\textbf{2. 特征根有重根的情形}

设特征方程有 $k$ 重根 $\lambda=\lambda_1$，则
\[
F(\lambda_1)=F'(\lambda_1)=\cdots=F^{(k-1)}(\lambda_1)=0,\qquad
F^{(k)}(\lambda_1)\neq0.
\]

(1) 若 $\lambda_1=0$，则特征方程含因子 $\lambda^k$，从而
$a_n=a_{n-1}=\cdots=a_{n-k+1}=0$，特征方程形如
\[
\lambda^n+a_1\lambda^{n-1}+\cdots+a_{n-k}\lambda^k=0,
\]
对应的方程变为
\[
\frac{d^n x}{dt^n}+a_1\frac{d^{n-1}x}{dt^{n-1}}+\cdots
+a_{n-k}\frac{d^k x}{dt^k}=0.
\]
它显然有 $k$ 个线性无关的解 $1,t,t^2,\dots,t^{k-1}$.

(2) 若 $\lambda_1\neq0$，作变换 $x=ye^{\lambda_1t}$，代入 (\ref{equ:3.15}) 并整理，
得
\[
L\left[ye^{\lambda_1t}\right]
=\left(\frac{d^n y}{dt^n}+b_1\frac{d^{n-1}y}{dt^{n-1}}+\cdots+b_ny\right)
e^{\lambda_1t}=L_1[y]e^{\lambda_1t},
\]
其中
\[
L_1[y]=\frac{d^n y}{dt^n}+b_1\frac{d^{n-1}y}{dt^{n-1}}+\cdots
+b_{n-1}\frac{dy}{dt}+b_ny,
\]
$b_1,b_2,\dots,b_n$ 仍为常数. 方程 $L_1[y]=0$ 的特征方程为
\[
G(\mu)=\mu^n+b_1\mu^{n-1}+\cdots+b_{n-1}\mu+b_n=0.
\]
由上述变换可得 $F(\mu+\lambda_1)=G(\mu)$. 因此 $F$ 与 $G$ 的导数也有相同的平移
关系，即 $\lambda_1$ 是 $F$ 的 $k$ 重根当且仅当 $0$ 是 $G$ 的 $k$ 重根. 这样就把
问题化为已讨论过的零根情形.

于是，$G(\mu)=0$ 的 $k$ 重零根对应 $L_1[y]=0$ 的 $k$ 个解
$1,t,t^2,\dots,t^{k-1}$，从而 (\ref{equ:3.15}) 有 $k$ 个解
\[
e^{\lambda_1t},\ te^{\lambda_1t},\ t^2e^{\lambda_1t},\ \dots,\
t^{k-1}e^{\lambda_1t}.
\]

一般地，设 (\ref{equ:3.17}) 的全部互异根为
$\lambda_1,\lambda_2,\dots,\lambda_m$，重数依次为 $k_1,k_2,\dots,k_m$，且
$k_1+k_2+\cdots+k_m=n$，则 (\ref{equ:3.15}) 有 $n$ 个解
\[
\left\{
\begin{aligned}
&e^{\lambda_1t},\ te^{\lambda_1t},\ t^2e^{\lambda_1t},\ \dots,\
t^{k_1-1}e^{\lambda_1t},\\
&e^{\lambda_2t},\ te^{\lambda_2t},\ t^2e^{\lambda_2t},\ \dots,\
t^{k_2-1}e^{\lambda_2t},\\
&\qquad\vdots\\
&e^{\lambda_mt},\ te^{\lambda_mt},\ t^2e^{\lambda_mt},\ \dots,\
t^{k_m-1}e^{\lambda_mt}.
\end{aligned}
\right.
\]

\begin{pf1}
证明这 $n$ 个解线性无关. 假设它们线性相关，则存在不全为零的常数
$C_j^{(r)}$，使得
\[
\sum_{r=1}^m\left[C_0^{(r)}+C_1^{(r)}t+\cdots
+C_{k_r-1}^{(r)}t^{k_r-1}\right]e^{\lambda_rt}
=\sum_{r=1}^mP_r(t)e^{\lambda_rt}\equiv0,
\]
其中 $P_r(t)$ 是次数不超过 $k_r-1$ 的多项式. 不妨设
$P_m(t)\not\equiv0$. 将上式除以 $e^{\lambda_1t}$，再对 $t$ 求导 $k_1$ 次，得
\[
\sum_{r=2}^mQ_r(t)e^{(\lambda_r-\lambda_1)t}\equiv0,
\]
其中 $Q_r(t)\equiv(\lambda_r-\lambda_1)^{k_1}P_r(t)+S_r(t)$，$S_r(t)$ 是次数低于
$P_r(t)$ 的多项式. 因此 $Q_m(t)$ 与 $P_m(t)$ 次数相同且 $Q_m(t)\not\equiv0$.
对上式重复同样的做法，经过 $m-1$ 次后得到
\[
R_m(t)e^{(\lambda_m-\lambda_{m-1})t}\equiv0,
\]
其中 $R_m(t)$ 与 $P_m(t)$ 次数相同且 $R_m(t)\not\equiv0$，矛盾. 故这 $n$ 个解
线性无关.
\end{pf1}

若特征方程有 $k$ 重复根 $\lambda=\alpha+\mathrm{i}\beta$，则
$\bar\lambda=\alpha-\mathrm{i}\beta$ 也是 $k$ 重复根. 与单复根的情形一样，可将
$2k$ 个复值解换成 $2k$ 个实值解
\[
\left\{
\begin{aligned}
&e^{\alpha t}\cos\beta t,\ te^{\alpha t}\cos\beta t,\ \dots,\
t^{k-1}e^{\alpha t}\cos\beta t,\\
&e^{\alpha t}\sin\beta t,\ te^{\alpha t}\sin\beta t,\ \dots,\
t^{k-1}e^{\alpha t}\sin\beta t.
\end{aligned}
\right.
\]

将上述结果归纳为如下步骤：

第一步 写出 (\ref{equ:3.15}) 的特征方程 (\ref{equ:3.17})，求出全部特征根.

第二步 按特征根的情形写出相应的解：
\begin{enumerate}
\item 对每个单实根 $\lambda_k$，有解 $e^{\lambda_kt}$；
\item 对每个 $m>1$ 重实根 $\lambda_k$，有 $m$ 个解
\[
e^{\lambda_kt},\ te^{\lambda_kt},\ t^2e^{\lambda_kt},\ \dots,\
t^{m-1}e^{\lambda_kt};
\]
\item 对每个单重共轭复根 $\alpha\pm\mathrm{i}\beta$，有两个解
\[
e^{\alpha t}\cos\beta t,\qquad e^{\alpha t}\sin\beta t;
\]
\item 对每个 $m>1$ 重共轭复根 $\alpha\pm\mathrm{i}\beta$，有 $2m$ 个解
\[
\left\{
\begin{aligned}
&e^{\alpha t}\cos\beta t,\ te^{\alpha t}\cos\beta t,\ \dots,\
t^{m-1}e^{\alpha t}\cos\beta t,\\
&e^{\alpha t}\sin\beta t,\ te^{\alpha t}\sin\beta t,\ \dots,\
t^{m-1}e^{\alpha t}\sin\beta t.
\end{aligned}
\right.
\]
\end{enumerate}

第三步 由第二步写出的 $n$ 个线性无关解构成基本解组，其线性组合即为
(\ref{equ:3.15}) 的通解.

\subsection{欧拉方程}
对一般的变系数齐次线性方程，没有普遍的代数解法. 但部分特殊类型可以通过自变量
替换化为常系数方程，欧拉方程就是典型代表.

\begin{deff}[(欧拉方程)]
形如
\[
t^n x^{(n)}+a_1t^{n-1}x^{(n-1)}+\cdots+a_{n-1}tx'+a_nx=0,\quad t>0
\]
的微分方程称为 \textbf{欧拉方程}，其中 $a_1,a_2,\dots,a_n$ 为实常数.
\end{deff}

作自变量替换 $t=e^s$，即 $s=\ln t$. 引入算子记号
\[
t\frac{d}{dt}=\frac{d}{ds},\qquad
t^2\frac{d^2}{dt^2}=\frac{d}{ds}\left(\frac{d}{ds}-1\right),
\]
更高阶有
\[
t^k\frac{d^k}{dt^k}
=\frac{d}{ds}\left(\frac{d}{ds}-1\right)\cdots
\left(\frac{d}{ds}-k+1\right).
\]

\begin{thm}\label{thm:3.3.2}
欧拉方程在变换 $t=e^s$ 之下化为关于未知函数 $x(s)$ 的常系数齐次线性方程.
\end{thm}
\begin{pf1}
将上述算子代入欧拉方程，即消去变量 $t$，得到只含 $s$ 的常系数齐次线性方程.
解出 $x(s)$ 后再作回代 $s=\ln t$，便得到原方程关于 $t$ 的解. 对 $t<0$ 的情形，
可作替换 $t=-e^s$ 类似处理.
\end{pf1}

\begin{rmk}
对二阶变系数齐次线性方程，若已求得一个非零特解 $x_1(t)$，则可用 Liouville
公式求出与之线性无关的第二个解
\[
x_2(t)=x_1(t)\int\frac{1}{x_1^2(t)}
\exp\left(-\int a_1(t)\,dt\right)dt,
\]
从而得到该方程的基本解组.
\end{rmk}

\section{线性非齐次常系数方程的待定系数法}\label{sec:3.4}
由 3.2 节非齐次线性方程解的结构定理及 3.3 节的结果，求非齐次常系数方程的通解，
只需给出对应齐次方程的通解与非齐次方程的一个特解，再将二者相加. 3.2 节已经给出
了求特解的常数变易法，但它需要积分运算，计算往往比较繁琐. 本节针对几类特殊的
非齐次项介绍更为简便的方法——\textbf{待定系数法}. 它的特点是不需要积分，只用
代数运算即可求得特解.

\subsection{非齐次项为多项式的情形}
考虑常系数非齐次线性方程
\begin{equation}\label{equ:3.19}
L[x]=\frac{d^n x}{dt^n}+a_1\frac{d^{n-1}x}{dt^{n-1}}+\cdots
+a_{n-1}\frac{dx}{dt}+a_nx=f(t),
\end{equation}
其中 $a_1,a_2,\dots,a_n$ 为常数，$f(t)$ 连续. 为叙述方便，以二阶方程为例介绍
比较系数法，先考虑
\begin{equation}\label{equ:3.20}
L[x]=\frac{d^2x}{dt^2}+p\frac{dx}{dt}+qx
=b_0+b_1t+\cdots+b_nt^n,
\end{equation}
其中 $p,q,b_0,b_1,\dots,b_n$ 为常数.

注意到多项式的一阶、二阶导数仍是多项式，而 (\ref{equ:3.20}) 的右端是 $n$ 次
多项式，因此可设特解为
\begin{equation}\label{equ:3.21}
\varphi(t)=t^k\left(B_0+B_1t+B_2t^2+\cdots+B_nt^n\right),
\end{equation}
其中 $B_0,B_1,\dots,B_n$ 为待定常数. 将 (\ref{equ:3.21}) 代入 (\ref{equ:3.20})，
比较两端同次幂的系数，即可确定这些常数.

(\ref{equ:3.20}) 对应的齐次方程的特征方程为
\begin{equation}\label{equ:3.22}
\lambda^2+p\lambda+q=0.
\end{equation}
根据 $q,p$ 的不同情形，(\ref{equ:3.21}) 中 $k$ 的取值如下.

当 $q\neq0$ 时，零不是 (\ref{equ:3.22}) 的根，取 $k=0$. 此时比较系数得
\[
\begin{cases}
qB_n=b_n,\\
qB_{n-1}+npB_n=b_{n-1},\\
\quad\cdots\cdots\\
qB_0+pB_1+2B_2=b_0.
\end{cases}
\]
由于 $q\neq0$，可自上而下逐个确定 $B_n,B_{n-1},\dots,B_0$.

当 $q=0,p\neq0$ 时，零是 (\ref{equ:3.22}) 的单根，取 $k=1$. 此时比较系数得
\[
\begin{cases}
(n+1)pB_n=b_n,\\
n(n+1)B_n+npB_{n-1}=b_{n-1},\\
\quad\cdots\cdots\\
6B_2+2pB_1=b_1,\\
2B_1+pB_0=b_0.
\end{cases}
\]
由于 $p\neq0$，同样可逐个确定各待定常数.

当 $q=0,p=0$ 时，零是 (\ref{equ:3.22}) 的二重根，取 $k=2$. 此时
$\varphi''=b_0+b_1t+\cdots+b_nt^n$，直接积分两次得
\[
\varphi(t)=t^2\left[\frac{b_0}{1\cdot2}+\frac{b_1t}{2\cdot3}
+\cdots+\frac{b_nt^n}{(n+1)(n+2)}\right].
\]
归纳以上讨论，(\ref{equ:3.20}) 有特解
\[
\varphi(t)=
\begin{cases}
B_0+B_1t+\cdots+B_nt^n, & q\neq0,\\
t\left(B_0+B_1t+\cdots+B_nt^n\right), & q=0,\ p\neq0,\\
t^2\left(B_0+B_1t+\cdots+B_nt^n\right), & q=0,\ p=0,
\end{cases}
\]
其中 $B_0,B_1,\dots,B_n$ 由比较系数法唯一确定.

\subsection{非齐次项为多项式与指数函数之积的情形}
其次考虑
\begin{equation}\label{equ:3.23}
L[x]=\frac{d^2x}{dt^2}+p\frac{dx}{dt}+qx
=\left(b_0+b_1t+\cdots+b_nt^n\right)e^{\alpha t},
\end{equation}
其中 $p,q,\alpha,b_0,b_1,\dots,b_n$ 为常数. 与 (\ref{equ:3.20}) 相比，
(\ref{equ:3.23}) 的右端多了一个因子 $e^{\alpha t}$. 为利用已有结论，作变换
$x=e^{\alpha t}y$，代入 (\ref{equ:3.23}) 得
\[
\frac{d^2y}{dt^2}+(2\alpha+p)\frac{dy}{dt}
+\left(\alpha^2+p\alpha+q\right)y=b_0+b_1t+\cdots+b_nt^n.
\]
由上一小节的结果，它有形如
\[
\psi(t)=
\begin{cases}
B_0+B_1t+\cdots+B_nt^n, & \alpha^2+p\alpha+q\neq0,\\
t\left(B_0+B_1t+\cdots+B_nt^n\right),
& \alpha^2+p\alpha+q=0,\ 2\alpha+p\neq0,\\
t^2\left(B_0+B_1t+\cdots+B_nt^n\right),
& \alpha^2+p\alpha+q=0,\ 2\alpha+p=0
\end{cases}
\]
的特解. 注意 $\alpha$ 是否为 (\ref{equ:3.22}) 的根，恰由
$\alpha^2+p\alpha+q$ 是否为零决定，故 (\ref{equ:3.23}) 的特解可归纳为：
\begin{enumerate}
\item 当 $\alpha$ 不是 (\ref{equ:3.22}) 的根时，取
$\varphi(t)=\left(B_0+B_1t+\cdots+B_nt^n\right)e^{\alpha t}$；
\item 当 $\alpha$ 是 (\ref{equ:3.22}) 的单根时，取
$\varphi(t)=t\left(B_0+B_1t+\cdots+B_nt^n\right)e^{\alpha t}$；
\item 当 $\alpha$ 是 (\ref{equ:3.22}) 的重根时，取
$\varphi(t)=t^2\left(B_0+B_1t+\cdots+B_nt^n\right)e^{\alpha t}$.
\end{enumerate}

\subsection{非齐次项为多项式、指数函数与三角函数之积的情形}
最后考虑
\begin{equation}\label{equ:3.24}
L[x]=\frac{d^2x}{dt^2}+p\frac{dx}{dt}+qx
=\left[A(t)\cos\beta t+B(t)\sin\beta t\right]e^{\alpha t},
\end{equation}
其中 $\alpha,\beta,p,q$ 为常数，$A(t),B(t)$ 为实系数多项式，其中一个的次数为
$m$，另一个的次数不超过 $m$.

由 Euler 公式，(\ref{equ:3.24}) 的右端可写成
\[
\frac{A(t)-\mathrm{i}B(t)}{2}e^{(\alpha+\mathrm{i}\beta)t}
+\frac{A(t)+\mathrm{i}B(t)}{2}e^{(\alpha-\mathrm{i}\beta)t}.
\]
于是 (\ref{equ:3.24}) 化为
\[
\frac{d^2x}{dt^2}+p\frac{dx}{dt}+qx
=\frac{A(t)-\mathrm{i}B(t)}{2}e^{(\alpha+\mathrm{i}\beta)t}
+\frac{A(t)+\mathrm{i}B(t)}{2}e^{(\alpha-\mathrm{i}\beta)t}.
\]
由非齐次线性方程解的叠加原理（定理 \ref{thm:3.2.11}），分别求右端两项对应的
方程的特解，再相加即可. 由于两个右端互为共轭，而方程的系数为实常数，若 $x_1$
是其中一个方程的解，则 $\overline{x_1}$ 是另一个方程的解.

对(\ref{equ:3.23}) 的结果（其中 $\alpha$ 可取复数）分别应用于这两个方程，可知
(\ref{equ:3.24}) 的特解具有形式
\[
\varphi(t)=t^k\left[P(t)\cos\beta t+Q(t)\sin\beta t\right]e^{\alpha t},
\]
其中 $P(t),Q(t)$ 是次数不超过 $m$ 的实系数多项式. 当
$\alpha+\mathrm{i}\beta$ 不是 (\ref{equ:3.22}) 的根时，取 $k=0$；当
$\alpha+\mathrm{i}\beta$ 是 (\ref{equ:3.22}) 的单根时，取 $k=1$.
